%! End Behavior and Asymptotes — Unit 2, Lesson 2.5 — Exit Ticket
\documentclass[10pt]{article}
\usepackage{atda-article}
\usepackage{atda-boxes}
\newcommand{\TallMath}[1]{$\displaystyle #1\rule[-1.4em]{0pt}{3.2em}$}
\setlength{\workrowsep}{8pt}

\begin{document}

\pageheader{Unit 2, Lesson 2.5}{Exit Ticket}
% NAMESTRIP: no \namedateperiod here. The student writes their name once, on the
% cover this component is stapled behind. See references/conventions.md.

\vspace{0.15in}

\begin{notesbox}{One Warming Soda, and a Graph Put On the Line}
\small
Notes closed. Three items. A cold soda is taken out of the refrigerator into a room held at
$72^\circ$F. The graph shows $T(t)$, its temperature in degrees Fahrenheit, $t$ hours later. The
dashed line is a landmark, not part of the graph.

\begin{center}
\begin{tikzpicture}
\begin{axis}[
    width=10.2cm, height=4.6cm,
    xlabel={$t$ (hours out of the refrigerator)}, ylabel={$T$ (degrees F)},
    xmin=0, xmax=5, ymin=0, ymax=80,
    xtick={0,1,2,3,4,5}, ytick={0,20,40,60,80}, minor y tick num=4,
    grid=both, grid style={linegray!45}, minor grid style={linegray!30},
    axis lines=left,
    tick label style={font=\tiny}, label style={font=\scriptsize}]
  \addplot[linegray, thick, dashed, domain=0:5, samples=2]{72};
  \addplot[royal, very thick, domain=0:5, samples=120]{72-64*(0.5)^x};
\end{axis}
\end{tikzpicture}
\end{center}

\begin{enumerate}[leftmargin=1.6em, itemsep=7pt, topsep=4pt]

\item Describe the end behavior of $T$ as $t$ runs off to the right. Give the horizontal asymptote as
an \emph{equation}, and say in a sentence what that line means about the soda.

\writelines{3}

\item Does $T$ have an absolute maximum? Does it have an absolute minimum? For each, give a value
\emph{and} a location, or explain why there is none.

\writelines{2}

\item A classmate writes: \emph{``The soda reaches $72^\circ$F at $t=5$ hours --- you can see the
graph sitting on the dashed line there.''} The number $72$ and the line are both correct answers to
item 1. Explain what is wrong with the sentence anyway, and write the correct statement.

\writelines{3}

\end{enumerate}
\end{notesbox}

\end{document}
